\documentclass[french]{neoschool}

\title{Le thème Celestia}
\subtitle{v1.0.0}
\author{Razik Ikhlef}
\date{\today}

\begin{document}

\maketitle

\section{Polices}
\begin{description}
\item[mainface] Police principale pour le texte et les titres (\textcolor{titleColor}{Literata} par défaut)
\item[mainfaceoptions] Options directement passées à \textcolor{titleColor}{fontspec} pour la police principale
\item[sansface] Police sans empattements pour les éléments structurels (\textcolor{titleColor}{Inter} par défaut)
\item[sansfaceoptions] Options passées à \textcolor{titleColor}{fontspec} pour la police sans empattements
\item[monoface] Police à chasse fixe pour le code (\textcolor{titleColor}{Roboto Mono} par défaut)
\item[monofaceoptions] Options passées à \textcolor{titleColor}{fontspec} pour la police à chasse fixe
\item[allserif] Utilise la police principale avec empattements pour les mathématiques
\end{description}

\begin{code}{latex}
\usetheme[
    mainface=EB Garamond,
    mainfaceoptions={Scale=1.1},
    sansface=Montserrat,
    monoface=Fira Code,
    allserif
]{Celestia}
\end{code}

\section{Couleurs}
\begin{description}
\item[maincolor] Couleur principale utilisée pour les titres et accents (code \textcolor{titleColor}{LaTeX} svgname ou \textcolor{titleColor}{HTML} hexadécimal)
\item[accentcolor] Couleur secondaire pour les éléments spéciaux
\item[backgroundcolor] Couleur de fond des diapositives (\textcolor{titleColor}{F7F9FC} par défaut)
\item[codebackgroundcolor] Couleur de fond des blocs de code (\textcolor{titleColor}{F1F3F6} par défaut)
\item[mainblue] Couleur des blocs standards (\textcolor{titleColor}{045549} par défaut)
\item[maingreen] Couleur des blocs exemple (\textcolor{titleColor}{054924} par défaut)
\item[mainred] Couleur des blocs alerte (\textcolor{titleColor}{490445} par défaut)
\item[unicolor] Utilise la couleur principale pour tout le texte
\end{description}

\begin{code}{latex}
\usetheme[
    maincolor=045549,
    accentcolor=E63946,
    backgroundcolor=FAFAFA,
    codebackgroundcolor=F5F5F5
]{Celestia}
\end{code}

\section{Mise en page}
\begin{description}
\item[margin] Marge du contenu (\textcolor{titleColor}{2em} par défaut)
\item[frametitle] Style du titre (\textcolor{titleColor}{elegant}, \textcolor{titleColor}{plain}, \textcolor{titleColor}{centered})
\item[decorative] Active les éléments décoratifs
\item[nodecorative] Désactive les éléments décoratifs
\item[centeredtitle] Centre le titre sur la page de titre
\item[titleright] Aligne le titre à droite
\end{description}

\begin{code}{latex}
\usetheme[
    margin=1.5em,
    frametitle=centered,
    decorative,
    centeredtitle
]{Celestia}
\end{code}

\section{En-têtes}
\begin{description}
\item[headstyle] Famille de police pour les titres : \textcolor{titleColor}{rmfamily} (avec empattements) ou \textcolor{titleColor}{sffamily} (sans empattements)
\item[headshape] Style des caractères : \textcolor{titleColor}{sc} (petites capitales), \textcolor{titleColor}{it} (italique), \textcolor{titleColor}{normal}
\item[headweight] Graisse des titres : \textcolor{titleColor}{bfseries} (gras) ou \textcolor{titleColor}{mdseries} (normal)
\end{description}

\begin{code}{latex}
\usetheme[
    headstyle=sffamily,
    headshape=sc,
    headweight=bfseries
]{Celestia}
\end{code}

\section{Code}
\begin{description}
\item[codebox] Active l'encadrement du code avec \textcolor{titleColor}{tcolorbox} (\textcolor{titleColor}{true} par défaut)
\item[nocodebox] Désactive complètement l'encadrement \textcolor{titleColor}{tcolorbox} du code
\item[nocodeframe] Conserve \textcolor{titleColor}{tcolorbox} mais sans bordure visible
\end{description}

\begin{code}{latex}
\usetheme[
    nocodebox,
    nocodeframe
]{Celestia}
\end{code}

\section{Pied de page}
\begin{description}
\item[nofooter] Supprime entièrement le pied de page, sauf le numéro de diapositive
\item[quartercirclefooter] Affiche uniquement le numéro dans un quart de cercle en bas à droite
\item[fullbarfooter] Crée une barre complète avec auteur/titre/date et numéro dans un cercle
\end{description}

\begin{code}{latex}
\usetheme[
    quartercirclefooter
]{Celestia}
\end{code}

\section{Table des matières}
\begin{description}
\item[compacttoc] Réduit l'espacement vertical entre les entrées de la table
\item[twocolumntoc] Répartit automatiquement les sections sur deux colonnes équilibrées
\end{description}

\begin{code}{latex}
\usetheme[
    compacttoc,
    twocolumntoc
]{Celestia}
\end{code}

\section{Blocs}
\begin{description}
\item[block] Bloc standard pour le contenu normal
\item[exampleblock] Bloc pour les exemples
\item[alertblock] Bloc pour les alertes
\end{description}

\section{Style des blocs}
\begin{description}
\item[soberblock] Le titre adopte la couleur principale du bloc (mainblue/maingreen/mainred) sur le fond général du document, tandis que le corps garde un fond légèrement teinté (10\%)
\item[softblock] Le titre et le corps partagent le même fond légèrement teinté (10\%), avec le titre dans la couleur principale correspondante
\end{description}

\begin{code}{latex}
\usetheme[
    soberblock
]{Celestia}

% ou

\usetheme[
    softblock
]{Celestia}
\end{code}

\section{Pages d'emphase}
L'option \texttt{standout} transforme une diapositive en page d'emphase, idéale pour les moments clés de la présentation~: citations marquantes, chiffres essentiels, messages à retenir.

\end{document}
